% Section 1.2, from lectures/L01/appendix/b_uart_top.md, with the provided reset synchronizer
% typeset straight from hw/reset_sync.vhd.

\section{Designing \code{uart_top.vhd}}
\label{c1:sec:uart_top}\label{c1:app:b}
\code{uart_top} is the peripheral: the datapath (\code{baud_gen}, \code{uart_tx}, \code{uart_rx}),
the register bank (\code{uart_regs}), and the provided SPI transport (\code{spi_slave},
\code{spi_reg_bridge}) that lets an off-chip processor reach the registers. It is built
\textbf{first}, top-down, before any of those blocks exist. That is the point of starting here: the
top defines every block's interface, so when \code{baud_gen} arrives in \chapref{2} or
\code{uart_regs} in \chapref{5}, the ports each one must expose are already fixed by how
\code{uart_top} instantiates it.

What this chapter produces is the top's \emph{skeleton}: the entity, the provided
\code{reset_sync}, the provided transport, and the handful of signals those three need. It is
wiring, almost entirely: the pieces of real logic this file eventually holds, the TX feeder and one
type conversion, are written later, in the chapters that bring the signals they act on. The
datapath and register blocks do not exist yet; each is instantiated here as it is created across
Chapters~2 to~5, and the system testbench turns green once the last of them lands.

The signals follow the same rule. Rather than declaring names for modules that will not exist for
weeks, each chapter reads the entity of the block it is about to instantiate and declares exactly
the signals that block's ports imply. Two conventions keep it readable: a \textbf{\code{spi_}
prefix} for signals carrying SPI traffic between the two transport blocks, and an \textbf{\code{_s2}
suffix} for anything that has been through a two-flop synchronizer, which is why every submodule
expects \code{reset_s2_n}.

This section is the reasoning behind the file. The exercises in \secref{c1:sec:exercises} are where
it is actually typed, port by port and signal by signal.

\subsection{Interface}
\label{c1:sec:uart_top_interface}

\bookfigure{uart_top}{Module \code{uart_top}}{c1:fig:uart_top}

\begin{booktable}{@{}p{8.6em}p{2em}p{6em}L@{}}
  \thead{Port} & \thead{Dir} & \thead{Type} & \thead{Meaning} \\
  \midrule
  \code{clock} & in & \code{std_logic} & 50\,MHz system clock. \\
  \code{reset_n} & in & \code{std_logic} & Asynchronous reset, active low (a button or a
    supervisor). \\
  \code{sclk}, \code{mosi}, \code{ss} & in & \code{std_logic} & SPI from the processor. \\
  \code{rx} & in & \code{std_logic} & The UART serial input. \\
  \code{miso} & out & \code{std_logic} & SPI to the processor. \\
  \code{tx} & out & \code{std_logic} & The UART serial output. \\
\end{booktable}

The ports are declared all inputs first, then all outputs; \code{uart_top_tb} binds them
positionally, so the order above is the contract. Nothing in this course associates a port by
name, which means a transposed pair of same-type pins is invisible to the analyzer, and stays
invisible until the system testbench runs at the end of \chapref{5}.

\code{uart_top} is the only block that takes \code{reset_n}, the raw asynchronous reset, from
outside. The one block inside it that also takes \code{reset_n} is \code{reset_sync}, because
manufacturing a clean reset is its whole job; every other block takes \code{reset_s2_n}, the
synchronized version it produces.

\subsection{Skeleton behaviour (wiring, and one placeholder)}
\label{c1:sec:uart_top_behaviour}
The transport comes first, because it is what makes a register bus exist at all. \code{spi_slave}
and \code{spi_reg_bridge} are given to you, testbench-verified; their wire protocol is Part~3 of the
protocol (\secref{proto:part3}), and their port lists are shown in Exercise~\ref{c1:ex:2}. You
instantiate them; you do not write them. \code{spi_slave} is the byte engine: it shifts bytes in and
out on \code{sclk} and reports each completed one. \code{spi_reg_bridge} assembles those bytes into
the five-byte register transactions the protocol defines, drives \code{reg_addr}, \code{reg_wdata}
and \code{reg_write}, and reads back whatever \code{reg_rdata} presents to it. The bank that will
answer on that bus arrives in \chapref{5}; until then the bus exists with nothing behind it.

The \textbf{reset synchronizer} is provided, as \code{reset_sync}, and you instantiate it.
\code{reset_n} comes from off-chip, so it can be released at any instant relative to \code{clock}. A
release landing near a clock edge could let different flip-flops leave reset on different cycles,
and a design that starts half-reset can sit in a state its own logic never expects.
\code{reset_sync} asserts asynchronously, so the moment \code{reset_n} goes low everything resets
with no clock needed, and releases synchronously, two flip-flops later, so the release is lined up
with the clock. Its two internal flops stay inside it, so the only signal it adds to your top is
\code{reset_s2_n}. The whole module is short enough to print:

\vhdlfile{hw/reset_sync.vhd}

\noindent That output then feeds every submodule, and the blocks you write in Chapters~2 to~5 all
take it the same way the provided ones do: \code{reset_s2_n} goes in the process's sensitivity
list, so the assertion still needs no clock edge, while the \emph{release} is already lined up
because \code{reset_sync} lined it up. That is the house style throughout \file{hw/}, and it is what
makes an asynchronous assertion safe in the first place: the dangerous edge, the release, was dealt
with once, here.

Note that this job is not \code{sync}'s (\chapref{3}), even though that module is also two flops in
series. \code{sync} takes \code{reset_s2_n} as an \emph{input}, so it consumes the very signal
\code{reset_sync} produces, and a block that needs a clean reset cannot be the block that
manufactures one.

The \textbf{TX feeder} is written in \chapref{5}, once both signals it needs exist. \code{uart_tx}
(\chapref{2}) sends a byte when its \code{start} is pulsed, and the TX FIFO inside \code{uart_regs}
(\chapref{5}) holds the bytes waiting to go. The feeder is what connects them: \code{tx_load} is
FIFO-not-empty AND transmitter-not-busy, and \code{tx_pop} is that same signal, so the transmitter
latches the front byte and the FIFO advances in the same cycle. The instant the transmitter goes
busy, \code{tx_load} drops, so exactly one byte moves per idle transmitter, never two. The RX side
needs no feeder at all: the receiver's \code{valid} and \code{data_out} become \code{uart_regs}'
\code{rx_push} and \code{rx_byte}, and the bank pushes its own FIFO.

The other piece is a \textbf{type conversion}, written in \chapref{2} alongside \code{baud_gen}.
\code{baud_gen} takes its divider as a \code{natural range 1 to 65535}, but \code{BAUD_DIV} leaves
the register bank as a \code{std_logic_vector(15 downto 0)}, so \code{uart_top} converts it. That
conversion is guarded, because until \code{uart_regs} exists in \chapref{5} nothing drives the
vector, and \code{to_integer} on an all-\code{'U'} value warns and substitutes a zero of its own
choosing, which is a value that port's range does not accept. Picking a legal value yourself costs
one conditional expression, and this is the only place in \code{uart_top} where a
\code{std_logic_vector} becomes an integer.

The one thing this chapter does write is a \textbf{placeholder}: \code{reg_rdata} tied to zeros, so
the bridge reads a defined value rather than \code{'U'} while the register bank is missing. It is
deleted in \chapref{5}, when \code{uart_regs} starts driving that vector for real.

\subsection{What the testbench pins down}
\label{c1:sec:uart_top_tb}
\code{uart_top_tb} is the system testbench, and it does not run until the last block lands in
\chapref{5}. It drives the peripheral entirely over SPI, as the processor would, and loops \code{tx}
back into \code{rx}. A single byte written over SPI is therefore transmitted, received through the
loopback, pushed into the RX FIFO, and read back over SPI: an end-to-end check that the reset
synchronizer, the transport, the register bank, the FIFOs, the feeder and both datapath halves all
agree. It also writes \code{BAUD_DIV} and reads it back, and checks \code{STATUS}, so the register
plumbing is exercised alongside the datapath.

Its one line of instantiation is also why the entity looks the way it does. The testbench binds
\code{port map(clock, reset_n, sclk, mosi, ss, rx, miso, tx)}, and that order, not the drawing in
Figure~\ref{c1:fig:uart_top}, is what the port table records.

\subsection{Where it fits}
\label{c1:sec:uart_top_fits}
\code{uart_top} faces two directions. Outward, its eight pins are the FPGA boundary: \code{sclk},
\code{mosi}, \code{ss} and \code{miso} reach the ATmega328P on the other side of the wire, and the
driver built in Chapters~6 to~9 talks to exactly these ports; \code{rx} and \code{tx} are the
serial line itself. In \chapref{9} a board wrapper sits above it and maps those pins onto DE0-CV
package pins, which is the only thing that wrapper does.

Inward, it is the one file revisited in every VHDL chapter, and each visit adds a block, the
signals that block's entity implies, and any glue those signals now make possible. This chapter
leaves it with the entity, \code{reset_sync}, the transport and a register bus nothing answers.
\chapref{2} adds \code{baud_gen} and \code{uart_tx}, and with them the \code{BAUD_DIV} conversion,
so the top can transmit. \chapref{4} adds \code{sync} on the \code{rx} pin, producing \code{rx_s2},
and \code{uart_rx} behind it, so the receiver never sees the raw asynchronous input. \chapref{5}
adds \code{uart_regs}, which owns both FIFOs, and the TX feeder that connects it to the transmitter;
with every block present, the system testbench elaborates for the first time. The rule that keeps
it compiling throughout is that you only ever instantiate an entity after you have created it.

\subsection{What's ahead}
The exercises in \secref{c1:sec:exercises} are, for \code{uart_top}, also the build sheet:
everything needed to type the file, from the entity down to the last internal signal, is there. In
\chapref{2} you build the first two datapath blocks, \code{baud_gen} and \code{uart_tx}, and
instantiate them here.
